\section{Limitations and open problems}
\label{sec:limitations}

\paragraph{What is proved, and in what domain.} All \st{UNSAT\_BOUNDED} results
are statements about a declared finite domain, under the declared Kähler chamber
$T_+$ and the declared target. The slice results remove the charge bound only on
slices satisfying \eqref{eq:star}. The only statement with no domain qualifier
is the divisibility obstruction. Empirical saturation is not a theorem.

\paragraph{Generic members.} All cohomology is that of a Zariski-generic member
of the relevant family---the full family for upstairs statements, the
$\Gam$-invariant family for downstairs ones. Cohomology can jump on special
members, which can add vector-like pairs or singlets. The leading-order Yukawa
matrices are an exception: every coupling evaluated uses Koszul identifications
valid for every smooth member.

\paragraph{The Yukawa engine.} Triple products are computed for hypersurfaces in
products of $\PP^1$ and for bundles whose $H^1$ sits at a single Koszul level.
Couplings with Koszul level sum two or more, which need a lift through the
defining polynomial, and classes spread over both Koszul levels are returned as
\st{UNKNOWN}. Neither occurs among the Yukawa couplings of any model studied;
both occur for some one-singlet couplings $SLH_u$ of published models, which
are therefore reported as undecided. Couplings with singlet insertions are higher products, not triple
products, and are decided only by the selection rule. The normalisation of the
trace map and the matter Kähler metrics (for which a localisation approach
exists~\cite{BBCLP2018}) are not computed, so no physical Yukawa coupling or
fermion mass is claimed.

\paragraph{Scope of the group theory.} Only abelian $\Gam$ acting without
permuting projective factors is handled; non-abelian groups are refused. The
free action and invariant family are input from the published symmetry data.

\paragraph{Published-data discrepancies left open.} 954 of 2,007 published down-type Yukawa
tables are internally inconsistent as charge tables (\cref{sec:yukawa}); we
report this without reconciling it. (A second apparent discrepancy, in 694
character tables, turned out to be an error in our own comparison convention and
is resolved in \cref{sec:equivariant}.)

\paragraph{Not addressed.} Supersymmetry breaking, moduli stabilisation, the
full $F$- and $D$-flatness of singlet vacua away from vanishing FI terms, and
anything non-perturbative. Rules R2 and R3 restrict the domain to that of
\cite{ACGLP2013}; searches for models with anti-families or positive pair indices
would need them relaxed, which the replays show is supported but costly.

\paragraph{Open problems.}
\begin{enumerate}[leftmargin=*]
\item \emph{Finiteness on the tetraquadric.} Prove a charge bound on the slices
  where $F$ is indefinite, by bounding the semialgebraic set
  $\{F\le0,\ \ind\in[-3|\Gam|,0]\}$ using the cubic growth of the index along
  $F\le0$ rays; equivalently, settle \cref{conj:C1}.
\item \emph{Higher Koszul levels.} Extend the triple-product evaluation to
  level sum two, and to complete intersections of higher codimension, following
  the lifting construction of~\cite{BBCL2016,BBL2017}.
\item \emph{Singlet vacua.} Decide which singlet vacuum values are $F$- and
  $D$-flat away from the Abelian locus, and hence which of the
  \st{FULL\_RANK\_ON\_LOCUS} verdicts can be realised.
\item \emph{Rank below $h-1$.} Decide exactly the 18 extra survivors whose charge
  matrix has rank below $h-1$, where the slope-zero locus is positive-dimensional
  and \cref{prop:decision} does not apply; and push the survey of
  \cref{sec:survey} to larger $K$, where more models missing from the catalog may
  appear.
\end{enumerate}
